\documentclass[11pt]{article}

\usepackage[utf8]{inputenc}

%\usepackage{fontspec}
%\usepackage[polish]{babel}

\usepackage{amsthm, amssymb}
\theoremstyle{definition}
\newtheorem{definition}{Definition}[section]
\newtheorem{theorem}{Theorem}[section]

\newtheorem{exercise}{Zadanie}[section]

% Technical.
\renewcommand{\phi}{\varphi}
\newcommand{\txt}[1]{\texttt{#1}}
\newcommand{\text}[1]{\texttt{#1}}
\renewcommand{\(}{\left(}
\renewcommand{\)}{\right)}

% Definitions.
\newcommand{\defn}{:\equiv}
\newcommand{\defp}{:=}

% Universes.
\newcommand{\U}{\mathcal{U}}

\newcommand{\isContr}{\text{isContr}}
\newcommand{\isProp}{\text{isProp}}
\newcommand{\isSet}{\text{isSet}}
\newcommand{\isGrpd}{\text{isGrpd}}

\newcommand{\Contr}{\text{Contr}}
\newcommand{\Prop}{\text{Prop}}
\newcommand{\Set}{\text{Set}}
\newcommand{\Grpd}{\text{Grpd}}
%\newcommand{\Type}{\text{Type}}

% Functions.
\newcommand{\lam}[2]{\lambda #1.#2}
\newcommand{\apl}[2]{#1\ #2}

\newcommand{\id}{\text{id}}
\newcommand{\comp}{\circ}
\newcommand{\finv}[1]{#1^{-1}}

% Paths.
\newcommand{\eq}[2]{#1 = #2}
%\newcommand{\refl}[1]{\text{refl}_{#1}}
\newcommand{\refl}{\txt{refl}}

\newcommand{\inv}[1]{#1^{-1}}
\newcommand{\sq}{\mathbin{\vcenter{\hbox{\rule{.3ex}{.3ex}}}}}

% Equivalence.
\newcommand{\hequiv}[2]{#1 \simeq #2}

% Path application and transport.
\newcommand{\ap}[2]{\text{ap}_{#1}(#2)}
\newcommand{\apd}[2]{\text{apd}_{#1}(#2)}
\newcommand{\transport}{\text{transport}}

\newcommand{\qinv}{\text{qinv}}
\newcommand{\isequiv}{\text{isequiv}}

% Basic types.
\newcommand{\Empty}{\mathbf{0}}
\newcommand{\Unit}{\mathbf{1}}

\newcommand{\Bool}{\mathbf{2}}
\newcommand{\true}{\txt{tt}}
\newcommand{\false}{\txt{ff}}

\newcommand{\Nat}{\mathbb{N}}
\newcommand{\Conat}{\mathbb{N}_\infty}
\newcommand{\dprod}[2]{\prod #1.#2}
\newcommand{\dsum}[2]{\sum #1.#2}

% Products.
\newcommand{\prodt}[2]{#1 \times #2}
\newcommand{\pair}[2]{(#1, #2)}
\newcommand{\outl}{\text{pr}_1}
\newcommand{\outr}{\text{pr}_2}

% Sums
\newcommand{\sumt}[2]{#1 + #2}
\newcommand{\inl}{\txt{inl}}
\newcommand{\inr}{\txt{inr}}

% Logic.
\newcommand{\all}[2]{\forall #1.#2}
\newcommand{\ex}[2]{\exists #1.#2}
\newcommand{\conj}[2]{#1 \land #2}
\newcommand{\disj}[2]{#1 \lor #2}
\newcommand{\True}{\top}
\newcommand{\False}{\bot}
\newcommand{\impl}[2]{#1 \implies #2}

% Rest.
\newcommand{\happly}{\text{happly}}
\newcommand{\funext}{\text{funext}}

\newcommand{\idtoeqv}{\text{idtoeqv}}
\newcommand{\ua}{\text{ua}}

\newcommand{\code}{\text{code}}
\newcommand{\encode}{\text{encode}}
\newcommand{\decode}{\text{decode}}

\newcommand{\hS}{\mathbb{S}^1}
\newcommand{\base}{\text{base}}
\newcommand{\looop}{\text{loop}}

\newcommand{\I}{\mathbb{I}}
\newcommand{\IZ}{0_\mathbb{I}}
\newcommand{\II}{1_\mathbb{I}}
\newcommand{\seg}{\text{seg}}

\newcommand{\trf}[1]{||#1||}
\newcommand{\tri}[1]{|#1|}

\title{Ścisłe zdania}

\begin{document}

\maketitle

\newcommand{\Stable}{\txt{Stable}}
\newcommand{\isStrictProp}{\txt{isStrictProp}}

\newcommand{\Decidable}{\txt{Decidable}}

\begin{center}
$\apl{\Stable}{X} \defn \neg\neg X \to X$ \\
\end{center}

\begin{exercise}
  Pokaż, że $\apl{\isProp}{(\apl{\Stable}{A})} \hequiv\ \apl{\isProp}{A}$
\end{exercise}

\newcommand{\dni}{\txt{dni}}

\begin{center}
$\dni : \dprod{A : \U}{A \to \neg \neg A}$ \\
$\apl{\apl{\dni}{a}}{na} \defn \apl{na}{a}$
\end{center}

\begin{exercise}
  Pokaż, że poniższe typy są równoważne:
  \begin{itemize}
    \item $\apl{\isContr}{(\apl{\Stable}{A})}$
    \item $\prodt{\apl{\isProp}{A}}{\apl{\Stable}{A}}$
    \item $\neg\neg A \to \apl{\isContr}{A}$
    \item $\apl{\isequiv}{\dni}$
    \item $A \hequiv\ \neg \neg A$
  \end{itemize}
\end{exercise}

\newcommand{\isStrictProp}{\txt{isStrictProp}}

\begin{center}
  $\apl{\isStrictProp}{A} \defn \prodt{\apl{isProp}{A}}{\apl{\Stable}{A}}$
\end{center}

\begin{exercise}
Które typy są stabilne, a które są stabilnymi zdaniami? Rozważ
\begin{itemize}
  \item $\Empty$
  \item $\Unit$
  \item $\U_i$
  \item Typy zamieszkałe
  \item $\prodt{A}{B}$
  \item $\sumt{A}{B}$
  \item $A \to B$
  \item $\neg{A}$
  \item $\dprod{x : A}{\apl{B}{x}}$
  \item $\dsum{x : A}{\apl{B}{x}}$
\end{itemize}

\end{exercise}

\newcommand{\StrictProp}{\txt{StrictProp}}

\begin{center}
  Ścisłe zdania

  \begin{itemize}
    \item $\StrictProp_i \defn \dsum{A : \U_i}{\apl{\isStrictProp}{A}}$
    \item $\bot \defn \Empty$
    \item $\top \defn \Unit$
    \item $P \land Q \defn \prodt{P}{Q}$
    \item $P \lor Q \defn \neg \neg (\sumt{P}{Q})$
    \item $P \Rightarrow Q \defn P \to Q$
    \item $P \Leftrightarrow Q \defn P = Q$
    \item $\neg P \defn P \to \Empty$
    \item $\all{x : A}{\apl{B}{x}} \defn \dprod{x : A}{\apl{B}{x}}$
    \item $\ex{x : A}{\apl{B}{x}} \defn \neg \neg \dsum{x : A}{\apl{B}{x}}$
  \end{itemize}
\end{center}

\begin{exercise}
  Pokaż, że ścisłe zdania mają klasykę:

  \begin{itemize}
    \item Umiarkowanie mocny wyłączony środek: $\all{A : \U}{A \lor \neg A}$
    \item Biedny wybór: $\all{A : \U}{\neg \neg (\neg \neg A \to A)}$
    \item Unikalny wybór: jeśli $P : X \to \U$ jest rodznią typów stabilnych, to $(\dprod{x : X}{\neg \neg \apl{P}{x}}) \to \dprod{x : X}{\apl{P}{x}}$
    \item Wybór jak 3.16 z książki: jeśli $X$ jest zbiorem, a $P$ jest rodziną ścisłych zdań, to $(\dprod{x : X}{\neg \neg \apl{P}{x}}) \hequiv\ \neg \neg \dprod{x : X}{\apl{P}{x}}$
    \item Skończony wybór: jeśli $X$ jest skończony, a $P$ dowolne, to \\ $(\dprod{x : X}{\neg \neg \apl{P}{x}}) \to \neg \neg \dprod{x : X}{\apl{P}{x}}$
  \end{itemize}
\end{exercise}

\newcommand{\acs}{\txt{acs}}

\begin{center}
Możemy przyjąć aksjomat wyboru: jeśli $X$ jest zbiorem, a $Y : X \to \U$ rodziną zbiorów, to mamy \\

$\acs : (\dprod{x : X}{\neg \neg \apl{Y}{x}}) \to \neg \neg \dprod{x : X}{\apl{Y}{x}}$

\end{center}

\begin{exercise}
  Pokaż, że trzeba uważać, żeby nie przedawkować klasyki:

  \begin{itemize}
    \item $\neg \dprod{A : \U}{\neg \neg A \to A}$
    \item Jeśli $X : \U$ dowolne, a $Y : X \to \U$ jest rodziną zbiorów, to \\ $\neg \left($(\dprod{x : X}{\neg \neg \apl{Y}{x}}) \to \neg \neg \dprod{x : X}{\apl{Y}{x}} \right)$
  \end{itemize}
\end{exercise}

\begin{exercise}
  Groźniejsze:

  \begin{itemize}
    \item Czy $\StrictProp$ ma resizing?
    \item Czy można udowodnić $\neg \neg \StrictProp \hequiv\ \txt{Bool}$?
  \end{itemize}
\end{exercise}

\begin{exercise}
  $\StrictProp$ jest dobrym kandydatem na uniwersum ścisłych zdań -- można dać jego typom wygodne reguły unikalności. Napisz te reguły.
\end{exercise}

\end{document}